\documentclass[conference]{IEEEtran}
\IEEEoverridecommandlockouts

\usepackage{cite}
\usepackage{amsmath,amssymb}
\usepackage{algorithmic}
\usepackage{graphicx}
\usepackage{textcomp}
\usepackage{xcolor}
\usepackage{booktabs}
\usepackage{url}
\usepackage{listings}
\usepackage[hidelinks]{hyperref}

\lstset{
  basicstyle=\ttfamily\scriptsize,
  columns=fullflexible,
  breaklines=true,
  frame=single,
  framesep=3pt,
  aboveskip=6pt,
  belowskip=4pt,
  keepspaces=true,
}

\def\BibTeX{{\rm B\kern-.05em{\sc i\kern-.025em b}\kern-.08em T\kern-.1667em\lower.7ex\hbox{E}\kern-.125emX}}

\begin{document}

\title{SixFields: Admission-Time Subtraction and Phase Folding\\for Cluster API}

\author{\IEEEauthorblockN{Madhava Gaikwad}
\IEEEauthorblockA{\textit{Independent} \\
yavan@outlook.com}}

\maketitle

\begin{abstract}
Cluster API is the standard Kubernetes mechanism for creating clusters
declaratively. In practice two things make it hard to use. Errors are reported
late and on objects the operator never wrote, and the wait during provisioning is
opaque, offering a single phase word with no estimate and no way to distinguish
slow from stuck. We describe SixFields, an assembly of upstream Cluster API that
attacks both without adding a custom resource or a controller. It restricts the
writable surface of a \texttt{Cluster} to six fields using a
\texttt{ValidatingAdmissionPolicy}, so a rejected field is reported at the moment
of the write and on the object that carries it; and it folds the condition trees
of roughly twenty generated objects into four phases with a percentile estimate
learned from prior runs, naming exactly one blocking object when progress stops.
The implementation separates a pure decision core from all cluster I/O through a
single snapshot envelope, so recorded runs replay through production code. We
report measurements from a single-operator deployment: 348 assertions in two
seconds, an admission policy exercised by 84 cases against both a CEL evaluator
and a live API server, and end-to-end cluster creation across four control-plane
configurations. Driving the tool against real clusters surfaced eleven defects
that fixtures had not, and eight behaviours of Cluster API v1.14.2 and k0smotron
v1.10.9 that are undocumented or contradict their documentation. We report those
findings, and we are explicit about what a single-machine, single-operator
evaluation cannot establish.
\end{abstract}

\begin{IEEEkeywords}
Kubernetes, Cluster API, admission control, developer experience, declarative
infrastructure, observability
\end{IEEEkeywords}

\section{Introduction}

Cluster API (CAPI) reconciles a declarative description of a Kubernetes cluster
into a running one \cite{capi}. An operator writes a \texttt{Cluster} object; a
set of controllers creates control-plane objects, machine objects, bootstrap
configurations and infrastructure objects, and drives them to readiness. The model
is sound and widely deployed. Using it is nevertheless difficult, and we argue the
difficulty concentrates in two places.

\textbf{Errors arrive late and elsewhere.} A \texttt{Cluster} that sets a field
owned by a controller is admitted without complaint. The field is later
overwritten, or it silently conflicts, and the visible symptom appears minutes
afterwards as a \texttt{False} condition on a \texttt{Machine} or an
infrastructure object the operator did not write and may not know exists. The
distance between cause and report, in both time and object identity, is the cost.

\textbf{The wait is opaque.} Provisioning takes minutes. During that window CAPI
exposes \texttt{status.phase} (for example \texttt{Provisioning}) and a tree of
conditions across every generated object. \texttt{clusterctl describe} renders the
tree \cite{clusterctl}, which is complete but undifferentiated: it does not say
which of the failing conditions is the cause, how far along the run is, or whether
the run has stopped making progress.

SixFields addresses these two, and deliberately nothing else. Our contributions:

\begin{enumerate}
\item \textbf{Subtraction as an admission-time property.} A CEL
\texttt{ValidatingAdmissionPolicy} \cite{vap} restricts \texttt{Cluster} writes to
six fields and refuses direct writes of the kinds a \texttt{ClusterClass}
generates. The rejection message is produced from the same template registry the
client uses, so a local pre-flight check and the server's rejection are
byte-identical (\S\ref{sec:subtraction}).
\item \textbf{A folding model for condition trees.} Four phases, a single ranked
blocking object, and a percentile estimate from local history
(\S\ref{sec:stream}). The ranking rules are stated, and each is motivated by a
misranking observed on a live cluster.
\item \textbf{An envelope indirection that makes replay equal production.} All
decision logic consumes one JSON snapshot and imports no Kubernetes client, so
recorded runs drive the same code paths as live ones (\S\ref{sec:envelope}).
\item \textbf{Empirical findings about CAPI v1.14.2 and k0smotron v1.10.9}
(\S\ref{sec:findings}), several of which contradict their own documentation, and
eleven defects that only appeared when the tool was driven against real clusters
(\S\ref{sec:defects}).
\end{enumerate}

\section{Background}

A \texttt{ClusterClass} is a template. It fixes what should be identical across
clusters, such as the infrastructure templates, the bootstrap provider and the
control-plane template, and exposes a small set of variables. A \texttt{Cluster}
with \texttt{spec.topology} names a class and supplies the differences. The
topology controller then generates the concrete objects.

CAPI v1.11 moved conditions to \texttt{metav1.Condition} under
\texttt{status.conditions}, retaining the previous generation under
\texttt{status.deprecated.v1beta1.conditions}; the contract for that generation is
\texttt{v1beta2} \cite{capicontract}. Conditions carry a type, a status, a reason,
a message and a transition timestamp. Some are positive-polarity, where
\texttt{True} is healthy. Others are negative-polarity, where the condition
reports an activity and \texttt{False} is the normal steady state. CAPI's own
aggregation code names the negative set explicitly.

Two properties of this model matter for what follows. First, a condition's
\texttt{lastTransitionTime} is the only signal available for detecting lack of
progress. Second, not every failure reaches a condition at all; we return to this
in \S\ref{sec:findings}.

\section{Design}

SixFields consists of three parts and adds no custom resource definition and no
controller. Everything is either declarative (a class, a policy, an add-on) or
client-side (one binary).

\subsection{Assembly}

One \texttt{ClusterClass}, \texttt{std}, with two user-facing variables:
\texttt{size} (\texttt{dev} or \texttt{ha}) and \texttt{placement} (\texttt{self}
or \texttt{hosted}). Three sibling classes exist: \texttt{std-hosted} for a hosted
control plane, \texttt{std-inmemory} for a simulated substrate, and
\texttt{std-kubeadm} as a documented fallback. A \texttt{ClusterResourceSet}
installs a digest-pinned CNI, selected by the
\texttt{topology.cluster.x-k8s.io/owned} label that CAPI places on every
class-derived \texttt{Cluster}, so no user field selects it.

\begin{figure}[t]
\begin{lstlisting}
spec:
  topology:
    classRef: {name: std}
    version: v1.34.11
    workers:
      machineDeployments:
      - {name: default, class: default, replicas: 2}
\end{lstlisting}
\caption{The writable surface. Six values, of which three are worker-pool
attributes. Two optional variables, \texttt{size} and \texttt{placement}, select
control-plane count and placement. Every other property of the cluster is fixed by
the class.}
\label{fig:six}
\end{figure}

Figure~\ref{fig:six} shows the whole of what an operator writes.

\subsection{Subtraction}
\label{sec:subtraction}

The writable surface is the object's name, namespace, labels and annotations;
\texttt{spec.topology.classRef}; \texttt{spec.topology.version}; each worker pool
entry's \texttt{name}, \texttt{class} and \texttt{replicas}; and variables
restricted to an allow-list. Everything else under \texttt{spec} is denied. A second policy denies
\texttt{CREATE} and \texttt{UPDATE} of the generated kinds by non-exempt
identities.

Three design points proved load-bearing.

\emph{Update semantics.} On \texttt{UPDATE} the policy compares \texttt{object}
against \texttt{oldObject} and denies only \emph{changes} to managed fields.
Without this, a controller writing \texttt{spec.controlPlaneRef} back onto the
\texttt{Cluster} would make every subsequent user edit fail.

\emph{Exemption.} Controller service accounts are exempted by username. An
ownership-label exemption is scoped to service-account requesters; unscoped it
would be a self-exemption hole, since \texttt{metadata.labels} is inside the
allowed surface.

\emph{Break-glass.} Bypass requires both a label on the object and membership of a
group. Either half alone is denied with a message naming the missing half, and
every use is recorded in the binding's audit annotation. Installing the assembly
is itself a break-glass write, because the class's templates are managed kinds.

Messages are rendered from one registry shared with the client, so
\texttt{cluster plan} reproduces the server's text exactly. A test asserts byte
equality between the two.

\subsection{Stream}
\label{sec:stream}

Folding maps the object set to four phases: infrastructure, control plane, workers
and add-ons. Each phase carries a state, one detail string, an elapsed time and
the set of objects contributing to it. Two rules emerged from live use. Only the
\emph{frontier} phase, the earliest that is not done, may be marked stalled;
workers sitting at $0/2$ behind a stuck control plane are waiting, not stuck.
Per-phase transition times are computed from only the conditions belonging to that
phase, because the \texttt{Cluster} object carries conditions for all four and
using its aggregate transition time would make every phase appear to change
whenever any did, defeating stall detection entirely.

\begin{figure}[t]
\begin{lstlisting}
Cluster/dev-1  class std . v1.34.11 . self
infrastructure  ############ done     ready              5s
control plane   ######______ running  0/1 nodes   ~13s left
workers         ######______ running  0/2 nodes
addons          ######______ running  waiting

safe to Ctrl-C; `cluster status dev-1` resumes
\end{lstlisting}
\caption{The ordinary display, mid-run. Four phases in the order they complete, a
state, one detail, and on the right either an elapsed time for a finished phase or
an estimate for the one the cluster is waiting on. Only the frontier phase carries
an estimate; the phases behind it are waiting, and their own history says nothing
about how long the phase in front will take. Reproduced from a recorded run;
block characters are substituted for the originals.}
\label{fig:phases}
\end{figure}

Figure~\ref{fig:phases} shows the display during a run.

\begin{figure}[t]
\begin{lstlisting}
DevMachine/dev-1-cp-abcde: etcd is not coming up (6m32s)
raw: kubectl get devmachine.infrastructure.cluster.x-k8s.io
     dev-1-cp-abcde -n default -o yaml
typical: p50 1m15s . p95 2m00s
next: cluster docs CAPI-CP-003
\end{lstlisting}
\caption{The stall block. One object, a plain-language cause, an elapsed time, the
command that produced the underlying object, the historical cost of this phase,
and a next action. The raw line is the only one permitted to exceed the terminal
width, because a truncated command is not copy-pasteable.}
\label{fig:stall}
\end{figure}

Figure~\ref{fig:stall} shows the output when a phase stalls.
Ranking selects one object when a phase stalls. Candidates are the \texttt{False}
conditions of contributing objects and of the objects they reference. Four rules
order them, in priority order:

\begin{enumerate}
\item \textbf{Specificity.} Leaf objects outrank their owners: infrastructure and
bootstrap objects, then \texttt{Machine}, then \texttt{MachineSet}, then pools,
then the control plane, then the \texttt{Cluster}.
\item \textbf{Cause over symptom.} A condition whose reason names another
condition on the same object that is also failing (for example
\texttt{WaitingForVMProvisioned}) ranks below the one it names.
\item \textbf{Specific over summary.} Aggregate conditions
(\texttt{Ready}, \texttt{Available}) rank below any specific condition on the same
object, since their text restates what a specific condition already reported.
\item \textbf{Recency, then severity.}
\end{enumerate}

Negative-polarity conditions are excluded from the candidate set entirely: they
report an activity, so \texttt{False} carries no information.

\begin{table}[t]
\caption{Ranking rules and the misranking each was introduced to fix. Every fix
was prompted by output observed on a live cluster.}
\label{tab:rank}
\centering
\footnotesize
\begin{tabular}{@{}p{0.30\columnwidth}p{0.60\columnwidth}@{}}
\toprule
Rule & Observed failure without it \\
\midrule
Exclude negative polarity & Reported \texttt{Paused=False}, meaning ``this object
is not paused'', as the blocking condition on a machine whose node had not come
up \\
Cause over symptom & Reported \texttt{APIServerProvisioned}, whose reason reads
\texttt{WaitingForVMProvisioned}, over the VM that was actually stuck \\
Specific over summary & Reported \texttt{Ready=False reason=NotReady}, which
restates, over \texttt{NodeHealthy} carrying the CNI error text \\
Specificity & Reported the \texttt{Cluster} over the \texttt{DevMachine} that
failed to pull an image \\
\bottomrule
\end{tabular}
\end{table}

Table~\ref{tab:rank} records each rule against the output that motivated it. We
note that three of the four were added after the tool was pointed at a real
cluster; the fourth was present from the design.

Estimates are percentiles over the last twenty completed runs of each phase, keyed
by (provider, class, placement, phase). With fewer than three samples the tool
reports that it has no history. No estimate is shown for a stalled phase; the
typical duration is shown instead, so the reader can judge how far off the run is.

\section{Implementation}

\subsection{The envelope and the pure core}
\label{sec:envelope}

All decision logic consumes a single JSON envelope holding every object behind one
cluster at one instant. The packages that fold, rank, estimate and render import
no Kubernetes client and take the clock as a parameter. One package,
\texttt{watch}, assembles the envelope from a live cluster.

The consequence is that a recorded run replays through production code. This is
the basis of the development loop (\texttt{cluster status --replay}), of the test
suite, and of the benchmark. It also removed an entire class of test double: there
are no mocks of the decision logic, because there is nothing to mock.

Recording is subject to redaction. A pre-release check found live kubeadm
bootstrap tokens in seven recorded fixtures; the recorder now redacts credentials
on write, and a test fails on any fixture containing a token-shaped string.

\subsection{Generated API truth}

CAPI's condition and reason identifiers change between releases. Rather than trust
recollection, a generator parses the pinned API modules out of the Go module cache
and emits a machine-readable snapshot: 659 constants with the file and line each
came from, plus a per-provider contract table. A test fails if any condition
identifier in the code is absent from that snapshot, and a documentation lint
fails if a user-facing page names one outside a code span.

\subsection{Executable interface contracts}

Every user-visible string is a typed entry in one registry. Because the registry
is enumerable, the interface contract is a test that walks it rather than a
convention. Enforced properties include: a next action on every message, by
construction, since the error constructor requires one; a stall line under 120
characters naming an object and an elapsed time; a denial naming the field, the
class and the break-glass; no raw condition identifiers in prose; a render budget
of at most two frames per second with no two consecutive identical frames; no
silence longer than the heartbeat while a run is in progress; and a versioned JSON
schema for machine consumers.

\subsection{Lints as interface enforcement}

Three properties are enforced by tests that walk the artefact rather than by
review. Condition identifiers in code must appear in the generated API snapshot,
which makes ``do not write an API detail from memory'' checkable. User-facing prose
may not name a condition identifier outside a code span, with the deny-list derived
from the same snapshot, so it updates itself when the upstream API does. And the
recorded fixtures may not contain credential-shaped strings.

We mention these because each replaced a convention that had already failed at
least once in the project's own history.

\section{Evaluation}

\subsection{Setup}

All measurements are from one machine: an Apple Silicon laptop with 24\,GB of
unified memory, running a kind management cluster under Docker Desktop. Pinned
versions: CAPI and CAPD v1.14.2 (contract \texttt{v1beta2}), k0smotron v1.10.9
(contract \texttt{v1beta1}), Kubernetes v1.34.11, k0s v1.34.11+k0s.0, Calico
v3.32.2. The optional model backend is \texttt{qwen2.5:14b} at Q4\_K\_M served by
Ollama.

The artefact is 7{,}137 lines of non-test Go, 4{,}018 lines of test Go, 1{,}280
lines of YAML and 2{,}981 lines of documentation, in 47 commits.

\subsection{Cost of the test suite}

\begin{table}[t]
\caption{Test layers, measured wall clock against declared budgets.}
\label{tab:tests}
\centering
\begin{tabular}{@{}llrr@{}}
\toprule
Layer & Substrate & Measured & Budget \\
\midrule
\texttt{make test} & none & 2\,s & 30\,s \\
\texttt{make test-envtest} & real API server & 16\,s & 180\,s \\
\texttt{make e2e} & simulated clusters & 187\,s & 600\,s \\
\bottomrule
\end{tabular}
\end{table}

Table~\ref{tab:tests} reports the three layers. The default layer runs 348
assertions across 140 test functions in two seconds with no cluster. Budgets are
asserted by the build, so a regression in test cost fails rather than accumulates.

The admission policy is exercised twice: 84 cases against a CEL evaluator reading
the policy YAML directly, and the same YAML applied to a real API server under
\texttt{envtest}, which type-checks every expression. A coverage test extracts the
denied key lists from the policy and fails if any lacks a negative case.

\subsection{End-to-end provisioning}

\begin{table}[t]
\caption{Time to \texttt{Available} by configuration, one run each.}
\label{tab:e2e}
\centering
\begin{tabular}{@{}llr@{}}
\toprule
Control plane & Substrate & Time \\
\midrule
kubeadm, machine-based & containers & 2\,m\,27\,s \\
k0smotron, hosted as pods & containers & 3\,m\,07\,s \\
k0s, machine-based & containers & 12\,m\,08\,s \\
kubeadm, machine-based & simulated & 44--60\,s \\
\bottomrule
\end{tabular}
\end{table}

Table~\ref{tab:e2e} gives time to \texttt{Available} for each configuration. These
are single runs on one machine and include image pulls in the container cases;
they establish that each path completes, and should not be read as a performance
comparison. The k0s machine-based path is slower because its bootstrap downloads
the k0s binary onto each machine at boot, a property we return to in
\S\ref{sec:findings}.

The escape hatch was verified on a live cluster: \texttt{cluster render} emitted
649 lines covering 18 objects, and \texttt{kubectl diff} reported no difference
against the running state.

\subsection{Ranking quality}

We evaluate ranking on six diagnosis tasks derived from stall fixtures, where
ground truth is the object and stall class the deterministic analyser assigns.
Three backends were compared: an identity backend that echoes the analyser, a
cassette replay, and a local language model. All three named the correct object
and class on all six tasks with no ungrounded output. Median latency was 250\,ns
for the analyser and 13\,s for the model.

We report this as a negative result for the plan the project set itself. The stated
rule was that a model-backed feature must beat the no-model baseline. On this
benchmark the baseline is also the ceiling: the analyser decides and the model only
explains, so the identity backend is correct by construction. What the benchmark
can measure is fidelity, meaning whether routing the analyser's finding through a
model preserves the object and the class, and the latency that costs. Every model
answer is checked against the input envelope; an answer naming an object or a
number absent from the envelope is discarded and the runbook stands alone.

\subsection{First-use study}

An agent with access only to the tutorial, the README and the binary, and
explicitly barred from reading source, attempted three tasks. It reached a correct
cluster status in 3 commands against a threshold of 3, and the cause of a stall in
1 command against a threshold of 2. It reported four documentation defects, all
since fixed. The most serious: three documents instructed the reader to verify the
escape hatch with \texttt{kubectl apply --dry-run=server}, which the admission
policy denies for thirteen of eighteen objects; and the managed-kind denial named
a class that did not own the object, because the policy cannot determine which of
four classes generated it.

This is a single trial with an automated subject and is not a user study. We
report it because it found defects that eleven days of self-review had not, and
because the tasks and thresholds were fixed before the run.

\subsection{Defects found only by live execution}
\label{sec:defects}

Eleven defects survived a fixture-backed test suite and appeared when the tool was
driven against real clusters. We group them because the pattern is the
contribution: each was invisible to synthetic data because the synthetic data
encoded the same misunderstanding as the code.

Four were in ranking or folding: negative-polarity conditions ranked as failures,
so the tool reported \texttt{Paused=False} (``this object is not paused'') as the
blocking condition; aggregate conditions outranking the specific ones carrying the
diagnosis; a symptom outranking the cause it named; and a
\texttt{ClusterResourceSetBinding} missing from the envelope, because it carries
no cluster label and is owned by the resource set, leaving the add-ons phase
reporting ``none'' on a cluster that had just had a CNI applied.

Four were in the client. Cobra's \texttt{Print} family writes to standard error, so
\texttt{cluster render > file} and \texttt{cluster new > file} both produced empty
files. One-shot commands opened a watch that never closed, so the command hung.
The replay clock was pinned to the synthetic epoch, so recorded fixtures replayed
with negative elapsed times. A one-shot status displaying a stall exited zero,
breaking the documented exit-code contract.

Three were in the watcher or the recorder: listing every served API version, so
each object was read twice with a deprecation warning per resource; going silent
once a cluster settled, so anything waiting on it waited forever; and writing
credentials into committed fixtures.

\section{Findings about the underlying systems}
\label{sec:findings}

The following behaviours of CAPI v1.14.2 and k0smotron v1.10.9 were established by
reading the pinned source and confirmed by execution. Several contradict project
documentation. We record them with citations to file and line in the artefact's
generated API snapshot.

\textbf{MachinePool publishes no v1beta2 conditions.} The constant block is
commented out in the pinned release with the note that it is not yet implemented.
Replica counters are populated. Any consumer folding a \texttt{MachinePool} must
read counters, not condition types.

\textbf{A MachineSet that cannot create a Machine reports it nowhere.} When
cloning a bootstrap template fails, for instance because a provider webhook rejects
it, \texttt{ScalingUp} remains \texttt{True} with the message ``Scaling up from 0
to 2 replicas'', no condition carries the reason, and no event is emitted. The
reason exists only in the controller log. This is the single largest gap we found
between what the API exposes and what an operator needs.

\textbf{Control-plane kind is immutable in three places.} A \texttt{Cluster}
cannot change placement after creation, cannot move to a class with a different
control-plane kind, and a \texttt{ClusterClass} cannot change its own
control-plane kind in place. Changing an assembly's bootstrap is therefore a
recreate, which is safe only while no cluster references the class.

\textbf{Template specs are immutable.} \texttt{Dev*Template.spec.template.spec}
cannot be updated, so two substrates require two template objects under two names.

\textbf{ClusterClass variables are defaulted onto the Cluster.} The mutating
webhook writes every class variable into \texttt{spec.topology.variables}. A
variable outside an admission allow-list therefore breaks an ordinary user's write.
Operator-controlled values must be written into templates, not exposed as
variables.

\textbf{Control-plane replicas cannot come from a class patch.} The
\texttt{KubeadmControlPlaneTemplate} has no replicas field and the patch engine
preserves \texttt{spec.replicas} on the control-plane object, so a patch is
discarded without error. The count must be written on the \texttt{Cluster}.

\textbf{k0smotron requires two spellings of one k0s release.}
\texttt{K0smotronControlPlane} takes the version as an image tag, where
\texttt{+} is invalid; \texttt{K0sWorkerConfig}'s webhook rejects the hyphenated
form and demands the \texttt{+k0s} suffix.

\textbf{K0sControlPlane reports a k0s version where CAPI expects a Kubernetes
version,} so CAPI's \texttt{ControlPlaneIsStable} preflight compares
\texttt{v1.34.11+k0s.0} against \texttt{v1.34.11}, concludes an upgrade is
permanently pending, and blocks worker scale-up indefinitely. Separately, a k0s
controller is not a Kubernetes node unless it also runs a worker, so its
\texttt{Machine} never obtains a node reference and never becomes ready.

\textbf{The in-memory backend's startup durations gate only at the first stage.}
Documentation and the project plan both suggest inducing a deterministic stall by
extending any component's \texttt{startupDuration}. Extending etcd's to thirty
minutes produced a cluster that reached readiness in 56 seconds. Only the VM stage
induces a stall, because its clock starts from a timestamp that always exists.

\section{Limitations}

The evaluation is a single operator on a single machine, and every timing is one
run. We make no claim about variance, about other hardware, or about clusters at
any realistic scale. The largest cluster built was three nodes.

The first-use study had one automated subject. It establishes that the documented
path is followable and that specific defects existed; it establishes nothing about
human users, learning time, or comparative usability against plain Cluster API. We
did not run the obvious control, which is the same tasks against
\texttt{clusterctl} alone, and without it the claim that the interface is an
improvement rests on the mechanism rather than on measurement.

The tool has never carried a production workload. The admission policy has been
exercised by 84 CEL cases, by a live API server, and by one operator; it has not
been subjected to an adversarial review.

Five of eleven fixtures remain synthetic, for stall classes that are awkward to
induce on demand. They are marked as such and a test enforces the marking, but
assertions resting on them rest on our model of the failure rather than on an
observed one.

Finally, the findings in \S\ref{sec:findings} are pinned to specific releases. They
are dated observations, and several are plainly defects that upstream may fix.

\section{Related work}

\texttt{clusterctl} \cite{clusterctl} is the reference client. Its
\texttt{describe} renders the full condition tree, and its
\texttt{alpha topology plan} performs a dry run of class changes. SixFields
consumes the same objects and differs in folding the tree to a fixed four phases,
ranking to a single object, and adding an estimate.

Giant Swarm and Syself both ship opinionated Cluster API platforms built on
\texttt{ClusterClass} and distributed as Helm charts, each targeting their own
customers' infrastructure. Their exposed value surfaces are the closest existing
analogue to the six fields here. Our differences are vendor neutrality, moving the
restriction into admission rather than leaving it to chart conventions, and
treating provisioning progress as a first-class interface.

Kubernetes \texttt{ValidatingAdmissionPolicy} \cite{vap} provides the CEL
admission mechanism; we contribute a use of it as a subtractive interface over an
existing API rather than as a conventional constraint checker. Policy engines such
as Gatekeeper and Kyverno could express the same rules; we chose the in-tree
mechanism to avoid adding a component to the cluster, which is a stated
non-goal of the design.

k8sgpt popularised running a deterministic analyser before any model and
anonymising what leaves the machine. We adopt both, and add a grounding check that
discards any model output naming an entity absent from the analyser's input.

\section{Conclusion}

Two properties of Cluster API's interface account for much of its difficulty:
errors that arrive late and on unfamiliar objects, and a provisioning wait that
exposes no progress. Both can be addressed without extending the API. Admission
policy moves a class of error to the moment and the object of the write. Folding a
condition tree into a fixed small number of phases, with one ranked blocking object
and an estimate from local history, makes the wait legible.

The engineering lesson we would carry forward is narrower and more transferable.
Routing all decision logic through a single recorded snapshot, so that replay and
production share code, was what made the system testable in two seconds. It was
also insufficient: eleven defects and eight upstream surprises appeared only when
the tool was pointed at a real cluster, because the fixtures encoded the same
assumptions as the code that consumed them. Recorded data is a strong substrate
for iteration and a weak substrate for discovery.

\section*{Availability}

Source, documentation and the generated API snapshot are at
\url{https://github.com/krimler/SixFields}, under Apache-2.0.

\begin{thebibliography}{00}
\bibitem{capi} Kubernetes SIG Cluster Lifecycle, ``The Cluster API Book,''
\url{https://cluster-api.sigs.k8s.io/}, accessed 2026.
\bibitem{clusterctl} Kubernetes SIG Cluster Lifecycle, ``clusterctl,''
\url{https://cluster-api.sigs.k8s.io/clusterctl/overview}, accessed 2026.
\bibitem{vap} Kubernetes, ``Validating Admission Policy,''
\url{https://kubernetes.io/docs/reference/access-authn-authz/validating-admission-policy/},
accessed 2026.
\bibitem{capicontract} Kubernetes SIG Cluster Lifecycle, ``Provider contracts,''
\url{https://cluster-api.sigs.k8s.io/developer/providers/contracts/overview},
accessed 2026.
\bibitem{cel} Google, ``Common Expression Language,''
\url{https://github.com/google/cel-spec}, accessed 2026.
\bibitem{k0smotron} k0smotron authors, ``k0smotron,''
\url{https://docs.k0smotron.io/}, accessed 2026.
\bibitem{kind} Kubernetes SIGs, ``kind: Kubernetes in Docker,''
\url{https://kind.sigs.k8s.io/}, accessed 2026.
\end{thebibliography}

\end{document}
